Teoria wymiaru jest jedną z dziedzin topologii, formalizujca tak zwane pojęcie „wymiaru“ przestrzeni topologicznej. Polega na przypisaniu każdej przestrzeni pewnej liczby, która zgodnie z naszą intuicją mówi, jak „duża“ jest dana przestrzeń. Najczęściej omawiane są mały wymiar indukcyjny, wymiar pokryciowy oraz duży wymiar indukcyjny.

Jednymi z najważniejszych osiągnięć teori wymiaru, jest pokazanie równości $\dim I^n=n$ \cite{bro1913} czego istotnym rezultatem jest pokazanie zależności $n\neq m \Longrightarrow \mathbb{R}^n\not \cong \mathbb{R}^m $, czyli że przestrzenie euklidesowe różnych wymiarów nie są homeomorficzne. W teorii wymiaru najczęściej rozważa się trzy funkcje: wymiar pokryciowy $\dim$, mały wymiar indukcyjny $\ind$ oraz duży wymiar indukcyjny $\Ind$. Dla przestrzeni metrycznych ośrodkowych zachodzi $ \dim X = \ind X = \Ind X.$ W ogólności równość ta nie jest spełniona.
Do ważnych osiągnięć należą też twierdzenia:
$\dim (X \cup Y) \le \dim X + \dim Y + 1$, czyli że suma dwóch przestrzeni zwiększa wymiar o co najwyżej jeden oraz $\dim (X \times Y) \le \dim X + \dim Y$, czyli wymiar iloczynu kartezjańskiego
nie przekracza sumy wymiarów.

Tematem tej pracy jest wymiar pokryciowy, który został poraz pierwszy wprowadzony przez Henri Lebesque'a \cite{lebesque1921}, a zformalizowany przez Eduarda Čecha \cite{cech1993}. 

W pierwszym rozdziale wprowadzimy definicję wymiaru pokryciowego, twierdzenie pozwalające nam zoptymalizować poszukiwanie go oraz wspomnimy o małym i dużym wymiarze powierzchniowym. W drugim rozdziale zdefiniujemy sympleksy oraz wielościany. Trzeci rozdział poświęcony będzie nerwom pokrycia, a w rodziale czwartym udowodnimy cel pracy, czyli twierdzenie o charakteryzacji przestrzeni topologicznych za pomocą $\varepsilon$-przekształceń. Głównym źródłem był podrzęcznik napisany przez Ryszarda Engelkinga \cite{engelking1995}.